% !TeX program = lualatex
% Compile twice: lualatex open_problems_500.tex
% All references and prose are embedded; no BibTeX database or external inputs.
\documentclass[11pt,a4paper]{article}
\usepackage[margin=24mm,headheight=14pt]{geometry}
\usepackage{fontspec}
\setmainfont{Latin Modern Roman}
\setsansfont{Latin Modern Sans}
\setmonofont{Latin Modern Mono}
\usepackage{amsmath,amssymb,mathtools}
\usepackage{unicode-math}
\setmathfont{Latin Modern Math}
\usepackage{microtype}
\usepackage{enumitem,xurl,needspace}
\usepackage[dvipsnames]{xcolor}
\usepackage{hyperref}
\hypersetup{unicode=true,colorlinks=true,linkcolor=MidnightBlue,urlcolor=MidnightBlue,
 pdftitle={Research notebook: Problems 40--49},pdfauthor={Source-annotated compilation},bookmarksnumbered=true}
\usepackage{bookmark}
\usepackage{tocloft}
\setlength{\cftsecnumwidth}{3em}
\cftsetpnumwidth{2.5em}
\cftsetrmarg{4em}
\usepackage{titlesec}
\titleformat{\section}{\Large\bfseries\raggedright}{\thesection}{0.7em}{}
\usepackage{fancyhdr}
\pagestyle{fancy}\fancyhf{}
\fancyhead[L]{\small\sffamily Mathematical problem statements}
\fancyhead[R]{\small Source edition 22}
\fancyfoot[C]{\thepage}
\setlength{\parindent}{0pt}
\setlength{\parskip}{5pt plus 1pt}
\setlength{\emergencystretch}{3em}
\setcounter{tocdepth}{1}
\setcounter{secnumdepth}{1}
\setlist[itemize]{leftmargin=3em,itemsep=5pt,topsep=4pt,parsep=0pt}
\allowdisplaybreaks
\newcommand{\N}{\mathbb{N}}
\newcommand{\Z}{\mathbb{Z}}
\newcommand{\Q}{\mathbb{Q}}
\newcommand{\R}{\mathbb{R}}
\newcommand{\C}{\mathbb{C}}
\newcommand{\F}{\mathbb{F}}
\newcommand{\A}{\mathbb{A}}
\newcommand{\PP}{\mathbb P}
\newcommand{\E}{\mathbb{E}}
\newcommand{\eps}{\varepsilon}
\newcommand{\dd}{\,\mathrm d}
\newcommand{\sref}[2]{\hyperlink{source:#1:#2}{[S#2]}}
\newcommand{\eref}[2]{\hyperlink{extra:#1:#2}{[E#2]}}
% Turn the next switch off for a shorter reading copy. The quotations preserve
% every exactTarget field, including qualifications omitted from a short summary.
\newif\ifshowcatalogue
\showcataloguetrue
\newcommand{\cataloguescope}[1]{\ifshowcatalogue
 \paragraph{Exact scope in the supplied catalogue.}
 \begin{quote}\small #1\end{quote}\fi}

% Research additions dated 27 September 2026; compile twice with LuaLaTeX.
\numberwithin{equation}{section}
\begin{document}
\setcounter{section}{45}
% ================================================================
% problemId: problem.fourier-restriction-conjecture
% source releaseRank: 46; source releaseStatus: open_with_solved_subcases
% exactTarget SHA-256: 979828e6e8c6d7ba269a150b9ab4c79e3deef889221d1d290d73f0f7013d9a54
\clearpage
\section{\texorpdfstring{Fourier restriction conjecture}{Fourier restriction conjecture}}\label{46}
\subsection{Definitions and mathematical statement}
For a Schwartz function on $\R^d$, set $\widehat f(\xi)=\int f(x)e^{-2\pi i x\cdot\xi}\,dx$, and let $d\sigma$ be surface measure on $S^{d-1}$. The restriction estimate is
\[
 \|\widehat f|_{S^{d-1}}\|_{L^q(d\sigma)}\le C_{d,p,q}\|f\|_{L^p(\R^d)}.
\]
In the usual range $1\le p,q<\infty$, the conjectured non-endpoint range is $p<2d/(d+1)$ and $q\le (d-1)p'/(d+1)$, with $p'=p/(p-1)$ and the usual interpretation at $p=1$. The full target asks for the sharp range in every $d\ge2$; endpoint conventions should be taken from the source, rather than inferred from a schematic inequality.
\par\smallskip\textit{Formulation and concept references:} \sref{46}{1}.
\cataloguescope{Prove the conjectured sharp L\textasciicircum{}p(R\textasciicircum{}d)-to-L\textasciicircum{}q(S\textasciicircum{}\{d-1\}) boundedness range for restriction of the Fourier transform to the unit sphere in every dimension d at least 2, including the non-endpoint range specified by the standard Knapp necessary conditions.}
\subsection{Short English statement}
Does Fourier restriction to a sphere satisfy every norm estimate allowed by the standard sharp necessary conditions?
% BEGIN RESEARCH ADDITION 46
\subsection{Research attempt}

\paragraph{Current frontier.}
\textit{Literature and formulation check, 27 September 2026.}
The source's Conjecture 3.5 uses a sphere of dimension one less than the
ambient dimension. After translating its notation, its strict radial
inequality and its non-strict sloping inequality agree with the statement
above. Its Theorem 3.7 gives the sharp $L^2$-restriction slice
$p\le 2(d+1)/(d+3)$; this is not the full conjectured range
\sref{46}{1}. The two catalogue entries for arXiv:2311.03145 refer to different
versions of one evolving manuscript. Version 15, dated 15 April 2026, concerns
a \emph{probabilistic analogue}; an estimate averaged over multipliers must
not be read as a bound for every deterministic multiplier
\eref{46}{1}. No resolution of the full stated range is inferred from these
sources. The unrestricted conjecture also has the Kakeya consequences
explained in the introduction of \sref{46}{1}.

\paragraph{Attack route.}
Write the extension operator as
\[
 Eg(x)=\int_{S^{d-1}} e^{2\pi i x\cdot\omega}g(\omega)\,d\sigma(\omega).
\]
Duality changes the target to $E:L^u(S^{d-1})\to L^r(\R^d)$ with
$u=q'$, $r=p'$. Three possible mechanisms are worth distinguishing.
Convolution geometry can control an even moment exactly; randomization can
remove cross terms but needs a deterministic replacement; and a multiscale
argument would have to bound the surviving resonances uniformly as the cap
size tends to zero. We pursue the first two, then test the proposed bridge
between them. The calculations below are derived here; the recovered
$L^2\to L^4$ estimate is a known special case, not a claimed improvement over
Stein--Tomas or a historical novelty claim.

\paragraph{Attempt.}
\textbf{1. A cap-local convolution lemma with an explicit constant.}
Suppose $d\ge3$, $0<c<1$, and $g$ is supported where
$\omega\cdot e\ge c$, for a unit vector $e$. For $0<|x|<2$, direct
integration on the sphere gives
\begin{equation}\label{eq:46:conv}
 (\sigma*\sigma)(x)=\frac{|S^{d-2}|}{|x|}
       \left(1-\frac{|x|^2}{4}\right)^{(d-3)/2}.
\end{equation}
For clarity about the normalization, set $x=\rho e_1$ and integrate
$\delta(|x-\omega|-1)$ against $d\sigma(\omega)$. In the coordinate
$v=\omega_1$, the spherical measure has density
$|S^{d-2}|(1-v^2)^{(d-3)/2}$; the root is $v=\rho/2$ and the absolute
value of the derivative of $|x-\omega|-1$ there is $\rho$. This proves
\eqref{eq:46:conv}. The density is zero outside the radius-two ball;
exceptional boundary values do not affect integration.

A sum of two points from the chosen cap has $x\cdot e\ge2c$ and hence
$|x|\ge2c$. Thus, on the support of the convolutions involving $g$,
\[
 (\sigma*\sigma)(x)\le K_{d,c}:=\frac{|S^{d-2}|}{2c}.
\]
Disintegrate $\sigma\otimes\sigma$ under $(\omega,\eta)\mapsto\omega+\eta$.
Cauchy--Schwarz on each fiber gives, almost everywhere,
\[
 |(g\sigma*g\sigma)(x)|^2
 \le (\sigma*\sigma)(x)
       \bigl((|g|^2\sigma)*(|g|^2\sigma)\bigr)(x).
\]
After integrating, the second convolution has total mass $\|g\|_2^4$.
Consequently,
\begin{equation}\label{eq:46:cap-four}
 \|Eg\|_4^4=\|g\sigma*g\sigma\|_2^2
       \le K_{d,c}\|g\|_2^4.
\end{equation}
Here the equality is justified rather than assumed: the inequality first
puts the convolution in $L^2$; its inverse Fourier transform is $(Eg)^2$
in distributions, so Plancherel identifies that distribution with an $L^2$
function. Since $g\sigma$ is a finite measure, $Eg$ is continuous, and this
identification also gives the asserted almost-everywhere equality of
functions. A finite cap cover and the $L^4$ triangle inequality yield a
global $L^2\to L^4$ bound. In $d=3$ this reaches the sharp $L^2$ slice;
in higher dimensions it is weaker than the quoted sharp slice.

\textbf{2. An exact test of the randomization step.}
For arbitrary complex numbers $z_1,\ldots,z_N$ and independent random signs
$\varepsilon_j$, expansion of the fourth power yields
\begin{align}
 \mathbb E\left|\sum_j\varepsilon_jz_j\right|^4
 &=2\left(\sum_j|z_j|^2\right)^2+
       \left|\sum_jz_j^2\right|^2-2\sum_j|z_j|^4\notag\\
 &\le3\left(\sum_j|z_j|^2\right)^2.
 \label{eq:46:random}
\end{align}
Indeed, a sign expectation vanishes unless every index occurs an even
number of times. The three pairings count each all-equal quadruple three
times, so two copies must be subtracted. This also proves the formula for
complex, rather than only real, $z_j$.

Taking every $z_j=1$ makes the deterministic fourth power $N^4$, whereas
the random expectation is $3N^2-2N$. There is therefore no dimension-free
pointwise replacement of the random average by the value at the all-positive
signs. For a finite cap decomposition $g=\sum_jg_j$, one may put
$z_j=Eg_j(x)$ and integrate \eqref{eq:46:random} over a finite ball by
Tonelli. That legal interchange does not supply the missing deterministic
inequality. Nonnegative cap data already have coherent phases at $x=0$.
This observation tests an intermediate inference, not the conjectured
whole-space norm bound.

\textbf{3. Strict curvature does not eliminate resonant rectangles.}
A stronger attempted shortcut would assert that four unit vectors satisfying
$\omega_1+\omega_2=\omega_3+\omega_4$ must coincide in pairs once all four
lie in a small cap. It is false. In $\R^3$, let
\[
 a=\delta e_1,\quad b=\delta e_2,\quad
 c_0=\sqrt{1-2\delta^2}\,e_3,\qquad 0<\delta<1/\sqrt2,
\]
and set
\[
 \omega_1=c_0+a+b,\quad \omega_2=c_0-a-b,\quad
 \omega_3=c_0+a-b,\quad \omega_4=c_0-a+b.
\]
All four vectors are distinct, have norm one, and have the same positive
third coordinate. Their two sums both equal $2c_0$. Letting $\delta\to0$
places them in arbitrarily small caps. The construction embeds in every
higher ambient dimension. Thus the fourth-moment expansion contains genuine
non-pairing resonances even on a strictly curved surface. Their removal by
independent signs is a change of problem; a deterministic proof must estimate
them, not declare them absent.

\textbf{4. Scaling identifies the range still missing.}
Choose $g_\delta$ as the indicator of a spherical cap of angular radius
$\delta$. After removing the common phase, phase oscillation is bounded by
a small constant on a box with $d-1$ tangential side lengths comparable to
$\delta^{-1}$ and one normal side length comparable to $\delta^{-2}$.
Taking the implicit constants small makes the real part of the integrand
positive, so
\[
 \|Eg_\delta\|_r\gtrsim
 \delta^{d-1-(d+1)/r},\qquad
 \|g_\delta\|_u\asymp\delta^{(d-1)/u}.
\]
This proves the necessary inequality
\begin{equation}\label{eq:46:knapp}
 r\ge\frac{d+1}{d-1}u'
\end{equation}
for finite $r$, with the usual interpretation when an exponent is infinite.
Smooth cap cutoffs give the same powers. These are necessary conditions,
not a proof that the conditions suffice.

For example, in $d=3$ and $u=\infty$, the target includes every $r>3$,
whereas \eqref{eq:46:cap-four} and the elementary $L^\infty$ bound give only
$r\ge4$ on the finite-measure sphere. Interpolation does not fill $3<r<4$.
On $\R^3$, the bounded function $(1+|x|)^{-\alpha}$ lies in $L^4$ but not
$L^r$ whenever $3/4<\alpha\le3/r$. This is not an extension-operator
counterexample; it falsifies the general integrability implication that the
interpolation shortcut would require.

\paragraph{Stress test.}
The convolution argument excludes $d=2$ deliberately: the exponent in
\eqref{eq:46:conv} would be negative and its singularity at $|x|=2$ is not
removed by the cap condition. No conclusion about the planar theorem is
being extracted from that argument. Antipodal singularities in $d\ge3$
are controlled locally before taking the finite cap sum. The calculation
neither assumes globally bounded $\sigma*\sigma$ nor interchanges an
infinite cap sum with an integral.

The random-sign identity, the resonant rectangle and the Knapp exponents
are exact. None bounds the coherent resonances at arbitrarily many scales.
A successful continuation needs a deterministic estimate below the obtained
$L^4$ exponent, with constants independent of cap resolution and spatial
radius, and with the input exponent linked to \eqref{eq:46:knapp}. The
probabilistic manuscript does not supply that missing universal quantifier
\eref{46}{1}. Finite sign enumeration and symbolic rectangle checks in the
accompanying script test the algebra only.

\paragraph{Outcome.}
\textbf{Outcome: Exploratory approach with unresolved gap.}
The attempt supplies a complete cap-convolution proof of a known special
case and two explicit obstructions to upgrading a randomized argument:
coherent fourth moments and non-pairing resonant rectangles within arbitrarily
small curved caps. These are counterexamples to proposed intermediate
shortcuts, not counterexamples to Fourier restriction.

The sharp all-dimensional range remains unproved here. In particular,
the argument does not produce the below-$L^4$ deterministic estimates already
needed in ambient dimension three, and it makes no claim to have resolved
the associated Kakeya implications or an endpoint outside the original scope.
% END RESEARCH ADDITION 46
\subsection{Sources}
\begin{itemize}\raggedright
\item[\textnormal{[S1]}]\hypertarget{source:46:1}{} Diogo Oliveira e Silva, Communications in Mathematics 32 (2024).
\url{https://cm.episciences.org/13778/pdf}.
{\small\textit{Catalogue formulation source.}}
\item[\textnormal{[S2]}]\hypertarget{source:46:2}{} Smooth Alpert frames, version 5.
\url{https://arxiv.org/abs/2311.03145v5}.
{\small\textit{Catalogue research/status reference; not a proof endorsement.}}
\item[\textnormal{[S3]}]\hypertarget{source:46:3}{} A probabilistic analogue of the Fourier extension conjecture.
\url{https://arxiv.org/abs/2311.03145}.
{\small\textit{Catalogue research/status reference; not a proof endorsement.}}
\item[\textnormal{[S4]}]\hypertarget{source:46:4}{} A survey of Stein's restriction conjecture.
\url{https://mathematics.stanford.edu/events/survey-steins-restriction-conjecture}.
{\small\textit{Catalogue research/status reference; not a proof endorsement.}}
% BEGIN RESEARCH SOURCES 46
\item[\textnormal{[E1]}]\hypertarget{extra:46:1}{} Eric T. Sawyer,
\emph{A probabilistic analogue of the Fourier extension conjecture},
arXiv:2311.03145v15, 15 April 2026; abstract and the distinction between
random multiplier estimates and deterministic extension estimates.
\url{https://arxiv.org/abs/2311.03145v15}.
{\small\textit{Version-specific research reference, consulted 27 September
2026. The older title in [S2] is retained unchanged; this entry is not an
endorsement of a deterministic resolution.}}
% END RESEARCH SOURCES 46
\end{itemize}

\end{document}
